\documentclass[11pt,a4paper]{article}
\usepackage[utf8]{inputenc}
\usepackage[T1]{fontenc}
\usepackage[spanish]{babel}
\usepackage[a4paper,margin=2.5cm]{geometry}
\usepackage{amsmath,amssymb,amsthm}
\usepackage{graphicx}
\usepackage{booktabs}
\usepackage{multirow}
\usepackage{tabularx}
\usepackage{longtable}
\usepackage{xcolor}
\usepackage{listings}
\usepackage{tikz}
\usetikzlibrary{arrows.meta,positioning}
\usepackage{hyperref}
\newcommand{\R}{\mathbb{R}}
\title{Documento de muestra}
\author{Equipo LaTeX visual}
\date{\today}
\begin{document}
\maketitle

\section{Ecuaciones}\label{sec:ecuaciones}
% align con cases dentro
\begin{align}
f(x) &= \begin{cases} x^2 & \text{si } x \ge 0 \\ -x & \text{si } x < 0 \end{cases} \label{eq:f}\\
g(x) &= \begin{pmatrix} 1 & 0 \\ 0 & 1 \end{pmatrix} x
\end{align}
\begin{gather}
a = b + c \\
\begin{bmatrix} a & b \\ c & d \end{bmatrix} \nonumber
\end{gather}
\begin{equation}
\begin{aligned}
x &= 1 \\ y &= 2
\end{aligned}
\label{eq:sistema}
\end{equation}
\begin{multline}
p(x) = a_0 + a_1 x + a_2 x^2 + a_3 x^3 \\
+ a_4 x^4 + a_5 x^5
\end{multline}
\begin{flalign}
u &= v & w &= z
\end{flalign}
\begin{eqnarray}
a & = & b + c \label{eq:arr1} \\
d & = & e \label{eq:arr2}
\end{eqnarray}
\begin{displaymath}
\int_0^1 x\,dx = \frac{1}{2}
\end{displaymath}
\begin{math}
e^{i\pi} + 1 = 0
\end{math}
\begin{alignat}{2}
x &= 1 &\quad y &= 2
\end{alignat}
Ver las ecuaciones \ref{eq:f}, \ref{eq:sistema}, \ref{eq:arr1} y \ref{eq:arr2}.

\section{Listas}
\begin{description}
\item[Teorema] Un enunciado demostrado.
\item[Lema] Un resultado auxiliar.
\end{description}
\begin{itemize}
\item[a)] Primera opción con etiqueta propia.
\item[] Sin viñeta.
\item Normal.
\end{itemize}

\section{Tablas}
\begin{table}[htbp]
\centering
\caption{Resultados agrupados}\label{tab:grupos}
\begin{tabular}{|l|c|c|}
\hline
\multirow{2}{*}{Grupo} & \multicolumn{2}{c|}{Medición} \\ \cline{2-3}
 & A & B \\ \hline
Uno & 1 & 2 \\ \hline
\end{tabular}
\end{table}
\begin{tabularx}{\linewidth}{lX}
\toprule
Campo & Descripción larga que se ajusta al ancho disponible \\
\bottomrule
\end{tabularx}
\begin{longtable}{ll}
Clave & Valor \\
uno & 1 \\
\end{longtable}
\begin{table*}[t]
\centering
\begin{tabular}{cc}
x & y \\
\end{tabular}
\end{table*}

\section{Figuras y espacios}
\begin{center}
\includegraphics[width=0.3\linewidth]{figura.png}
\end{center}
\begin{figure}[htbp]
\centering
\begin{minipage}{0.45\linewidth}
\centering
Texto de la izquierda con \textbf{negrita}.
\end{minipage}\hfill
\begin{minipage}{0.45\linewidth}
\centering
$E = mc^2$
\end{minipage}
\caption{Dos columnas con minipage}\label{fig:mini}
\end{figure}
\hrule
\vspace{1cm}
Texto en \textcolor{red}{rojo}, en \textsc{Versalitas}, código \verb|x_1 = 2| y \lstinline|print(x)|, fecha \today, el logo \LaTeX, {\small texto pequeño} y {\Large \bfseries grande y en negrita}. Un dibujo en línea: \tikz \draw[fill=blue] (0,0) circle (2pt);.

\section{Capítulo incluido}
Este texto viene de capitulo1.tex.


\section{Dibujos}
\begin{tikzpicture}
\draw[thick, ->] (0,0) -- (2,0) node[right] {$x$};
\draw[thick, ->] (0,0) -- (0,1.5) node[above] {$y$};
\end{tikzpicture}

\begin{figure}[htbp]
\centering
\begin{tikzpicture}[scale=0.8]
\foreach \x in {0,1,...,4} {
  \draw (\x,0) circle (0.3);
  \node at (\x,-0.7) {\x};
}
\end{tikzpicture}
\caption{Círculos con foreach}\label{fig:foreach}
\end{figure}

\begin{figure}[htbp]
\centering
\begin{tikzpicture}[node distance=1.5cm]
\node[draw, rounded corners] (a) {Inicio};
\node[draw, right=of a] (b) {Proceso};
\draw[-{Stealth}] (a) -- (b);
\end{tikzpicture}
\caption{Con librerías positioning y arrows.meta}\label{fig:flujo}
\end{figure}

\begin{center}
\resizebox{0.5\linewidth}{!}{%
\begin{tikzpicture}
\draw (0,0) rectangle (3,2);
\end{tikzpicture}%
}
\end{center}
Las figuras \ref{fig:mini}, \ref{fig:foreach} y \ref{fig:flujo}; la tabla \ref{tab:grupos}.

% Un caso realmente exótico: un entorno propio
\newenvironment{caja}{\begin{quote}\itshape}{\end{quote}}
\begin{caja}
Contenido en una caja propia.
\end{caja}

\bibliographystyle{plain}
\bibliography{referencias}

\end{document}
